% Source control vertices transcribed verbatim from the user's TikZ message.
% This file preserves the original Hobby path; separate Python source
% uses periodic cubic splines and changes P10 y by +0.07 to regularize
% its coincidence with P3. Do not interpret this reference drawing as a
% certified KPC 3D model or transfer its crossing IDs without runtime audit.
\documentclass[tikz,border=6pt]{standalone}
% ------- TikZ Guide Lines -------
\newcommand{\SSTGuidesPoints}[2]{% #1=basename (e.g. P), #2=last index
  \ifsstguides
    \foreach \i in {1,...,#2}{
      \fill[blue] (#1\i) circle (1.2pt);
      \node[blue,font=\scriptsize,above] at (#1\i) {\i};
    }
    \draw[gray!40, dashed]
    \foreach \i [remember=\i as \lasti (initially 1)] in {2,...,#2,1} { (#1\lasti)--(#1\i) };
  \fi
}
% ------- TikZ Guide Lines -------

\usepackage{tikz}
\usetikzlibrary{hobby,knots}
\usetikzlibrary{decorations.markings}

\begin{document}
\begin{tikzpicture}[use Hobby shortcut]
    \tikzset{knot diagram/every strand/.append style={draw=black}}
    \coordinate (P1) at (0.75, 3.25);
    \coordinate (P2) at (-1, 2.75);
    \coordinate (P3) at (-3, 1.75);
    \coordinate (P4) at (-2.75, -0.25);
    \coordinate (P5) at (-1.75, -2.50);
    \coordinate (P6) at (0.25, -1.75);
    \coordinate (P7) at (-1, -1.25);
    \coordinate (P8) at (-3.25, -0.50);
    \coordinate (P9) at (-2.50, 1.50);
    \coordinate (P10) at (-1.25, 3.25);
    \coordinate (P11) at (0.75, 2.75);
    \coordinate (P12) at (2.50, 1.75);
    \coordinate (P13) at (2, -0.25);
    \coordinate (P14) at (1, -2.50);
    \coordinate (P15) at (-0.75, -2);
    \coordinate (P16) at (0.50, -1.25);
    \coordinate (P17) at (2.75, -0.50);
    \coordinate (P18) at (2, 1.50);

    \begin{knot}[
        consider self intersections,
        clip width=5pt, clip radius=3pt,
        ignore endpoint intersections=false,
        flip crossing/.list={2,4,6,8,10,12,14,16,18,20},
    % ----draft mode=crossings % uncomment to see numbers
    ]
        \strand
        ([closed] P1)..(P2)..(P3)..(P4)..(P5)..(P6)..(P7)..(P8)..(P9)..(P10)..(P11)..(P12)..(P13)..(P14)..(P15)..(P16)..(P17)..(P18);
    \end{knot}
    \SSTGuidesPoints{P}{18}
\end{tikzpicture}
\end{document}